\usepackage{amssymb}% \varnothing
\usepackage{latexsym}

\usepackage[geometry]{ifsym}
\usepackage{amsmath}
\usepackage{amssymb}
\usepackage{latexsym}
\usepackage{relsize}
\usepackage{color,graphicx}
\usepackage{xfrac}


%\usepackage{boisik}
%\usepackage{lmodern}% http://ctan.org/pkg/lm
%\usepackage{amsmath}

\usepackage{stmaryrd}
%\SetSymbolFont{stmry}{bold}{U}{stmry}{m}{n}

\newtheorem{mydescription}{Descrição}[section]
\usepackage{calc}
\usepackage{enumitem}
\SetLabelAlign{parright}{\parbox[t]{\labelwidth}{\raggedleft#1}}

\usepackage{listings}
\usepackage{array,tabularx}
\usepackage{colortbl}
%\usepackage[table]{xcolor}

% FIGURE

\usepackage[font=scriptsize,skip=2pt]{caption}

\usepackage{tabularx} 
\usepackage{xcolor}
\newcommand{\cor}{\textcolor}
\usepackage{media9}
\usepackage{comment}

%LaTeX Font Warning: Font shape `OT1/cmr/m/n
%Timeline for How to remove the warnings “Font shape `
\usepackage{anyfontsize}
\usepackage{caption}
\usepackage{subcaption}
\usepackage{framed}

%\usepackage[lighttt]{lmodern}
\usepackage{listings}
\lstset{basicstyle=\ttfamily, keywordstyle=\bfseries}

\usepackage{multicol}
\usepackage{parskip}
\usepackage{pxfonts}
\usepackage{graphics}
\usepackage{pstricks}

% Texto justificado
\usepackage{ragged2e}
\usepackage{colortbl}
\usepackage{comment}

% ---------------------------------------------------
% Set font size for footnotes

%\makeatletter
%\renewcommand\footnotesize{%
%   \@setfontsize\scriptsize\@ixpt{7}%
%   \abovedisplayskip 8\p@ \@plus2\p@ \@minus4\p@
%   \abovedisplayshortskip \z@ \@plus\p@
%   \belowdisplayshortskip 4\p@ \@plus2\p@ \@minus2\p@
%   \def\@listi{\leftmargin\leftmargini
%               \topsep 4\p@ \@plus2\p@ \@minus2\p@
%               \parsep 2\p@ \@plus\p@ \@minus\p@
%               \itemsep \parsep}
%   \belowdisplayskip \abovedisplayskip
%}
%\makeatother

% Set font size for footnotes
%\renewcommand{\footnotesize}{\scriptsize}

%You can use arbitrary font sizes within LaTeX, not only the default ones 
% (\normalsize, \Large, \footnote size, &c.). The command for a font size of 9pt would be:

\renewcommand{\footnotesize}{\fontsize{6pt}{9pt}\selectfont}

%The second value is value for the interline skip for that size. 
%Usually it is about 20 % more than the nominal size.
% ---------------------------------------------------
%\AtBeginDocument{\fontsize{8}{10}\selectfont}

\UseRawInputEncoding
%\usepackage{ae,lmodern}
% Tabelas TcolorBox


\makeatother

%\usepackage{tcolorbox}
\tcbuselibrary{skins}
%\tcbset{beamer without segmentation/.style={%
%  beamer,segmentation code=,  interior titled code={{\tcb@spec{beamer@color}\tcb@drawwithtitle@path}\tcb@drawspec@T},  interior code={{\tcb@spec{beamer@color}\tcb@drawwithouttitle@path}}}}


\newcommand{\PR}[1]{\ensuremath{\left[#1\right]}}
\newcommand{\PC}[1]{\ensuremath{\left(#1\right)}}
\newcommand{\CH}[1]{\ensuremath{\left\{#1\right\}}}


%\newcommand\scaledfrac[2]{\scalebox{2}{\sfrac{#1}{#2}}}

%\newcommand{cmd}[args][opt]{def}
%\newtheorem{Exemplo}{Exemplo}[section]
% Closed” (square) root symbol

\usepackage{letltxmacro}
\makeatletter
\let\oldr@@t\r@@t
\def\r@@t#1#2{%
\setbox0=\hbox{$\oldr@@t#1{#2\,}$}\dimen0=\ht0
\advance\dimen0-0.2\ht0
\setbox2=\hbox{\vrule height\ht0 depth -\dimen0}%
{\box0\lower0.4pt\box2}}
\LetLtxMacro{\oldsqrt}{\sqrt}
\renewcommand*{\sqrt}[2][\ ]{\oldsqrt[#1]{#2}}
\makeatother


